\documentclass[11pt,a4paper]{article}
\usepackage[margin=2.5cm]{geometry}
\usepackage[expansion=false]{microtype}
\usepackage{amsmath,amssymb}
\usepackage{booktabs}
\usepackage{enumitem}
\usepackage{xcolor}
\usepackage{listings}
\usepackage[colorlinks=true,linkcolor=blue!50!black,citecolor=blue!50!black,urlcolor=blue!50!black]{hyperref}
\usepackage[capitalise,noabbrev]{cleveref}

\definecolor{codebg}{RGB}{246,246,246}
\definecolor{codecomment}{RGB}{0,110,0}
\lstset{language=Matlab, basicstyle=\ttfamily\small, backgroundcolor=\color{codebg},
  commentstyle=\color{codecomment}, keywordstyle=\color{blue!60!black}\bfseries,
  stringstyle=\color{purple!70!black}, frame=single, rulecolor=\color{gray!40},
  columns=fullflexible, keepspaces=true, showstringspaces=false, breaklines=true,
  xleftmargin=2pt, xrightmargin=2pt, aboveskip=6pt, belowskip=6pt}

\newcommand{\Sym}{\mathbb{S}}
\newcommand{\R}{\mathbb{R}}
\newcommand{\tr}{\operatorname{tr}}
\newcommand{\svec}{\operatorname{svec}}
\newcommand{\cA}{\mathcal{A}}
\newcommand{\cK}{\mathcal{K}}
\newcommand{\code}[1]{\texttt{#1}}
\setlength{\parskip}{3pt}
\emergencystretch=2em

\title{\textbf{SOSNOSDP}\\[4pt]\large Matrix sum-of-squares programs in SOSTOOLS,\\ solved without semidefinite programming\\[8pt]\normalsize User manual, version 0.1.0}
\author{Corentin Briat\thanks{www.briat.info; corentin@briat.info}}
\date{}
% October 2026}

\begin{document}
\maketitle

\begin{abstract}
\noindent SOSNOSDP is a MATLAB package (also running in GNU Octave) that solves sum-of-squares (SOS) programs written with SOSTOOLS without any semidefinite programming (SDP) solver. Constraint polynomials, scalar or matrix-valued, are represented by their values at interpolation nodes, and the program is solved by a non-symmetric interior-point method that only evaluates the barrier of the dual cone. This is the approach of Papp and Y{\i}ld{\i}z~\cite{PappYildiz2019}, in its matrix-valued form due to Kapelevich, Coey and Vielma~\cite{KCV2023} and implemented in the Julia solver Hypatia~\cite{CKV2022}. SOSNOSDP re-implements the matrix weighted SOS cone and the interior-point method in MATLAB, connects them to SOSTOOLS~\cite{SOSTOOLS}, adds interval constraints that need no multipliers, and certifies every returned solution. This manual explains how the package works and how to use it.
\end{abstract}

\tableofcontents

%======================================================================
\section{Introduction}\label{sec:intro}
%======================================================================
SOSTOOLS~\cite{SOSTOOLS} converts an SOS program into an SDP by writing every SOS polynomial with a positive semidefinite Gram matrix, and passes the SDP to a solver such as SeDuMi~\cite{Sturm1999}. The Gram matrices of a polynomial of degree $2d$ in one variable have about $d^2/2$ entries, while the polynomial has only $2d+1$ coefficients, and for an $m\times m$ matrix polynomial the Gram matrix has side $m(d+1)$. The lifting therefore dominates time and memory at high degree.

Papp and Y{\i}ld{\i}z~\cite{PappYildiz2019} showed that the lifting can be avoided entirely. A polynomial is stored by its values at unisolvent interpolation nodes, the SOS cone appears as the dual slack of a conic program, and the homogeneous interior-point method of Skajaa and Ye~\cite{SkajaaYe2015,PappYildiz2017} needs a barrier only for the dual of the SOS cone. That barrier, $-\log\det\Lambda(x)$, has cheap closed-form derivatives in the interpolant basis. Their MATLAB solver alfonso~\cite{alfonso} implements the scalar case. The extension to SOS matrices, using the SOS cones in general algebras of Papp and Alizadeh~\cite{PappAlizadeh2013}, is due to Kapelevich, Coey and Vielma~\cite{KCV2023}, and is implemented in Hypatia~\cite{CKV2022,CKV2023}.

SOSNOSDP brings this approach to SOSTOOLS users. Its contributions are practical:
\begin{itemize}[itemsep=1pt]
\item a MATLAB implementation of the matrix weighted SOS cone of~\cite{KCV2023} and of the homogeneous interior-point method;
\item \code{sossolve\_nosdp}, a replacement for \code{sossolve} that reads a SOSTOOLS program and solves it without SDP, after which \code{sosgetsol} works unchanged;
\item native interval constraints for scalar and matrix polynomials, without multipliers, and a helper for normalized time variables;
\item certificates for the coefficients actually returned, with explicit checks against rounding errors.
\end{itemize}
The theory of the cone, its barrier and the recovery of SOS decompositions is that of~\cite{PappYildiz2019,KCV2023}. Please cite those works together with this package (\cref{sec:cite}).

\subsection{Scope: one variable or several}\label{sec:scope}
Matrix SOS constraints are \emph{not} restricted to one variable. What depends on the number of variables is the kind of set on which positivity is imposed:
\begin{center}\small
\begin{tabular}{p{4.8cm}p{1.9cm}p{6.6cm}}
\toprule
Constraint & Variables & Treatment\\
\midrule
\code{sosineq}, \code{sosmatrixineq} (positivity on $\R^n$) & any number & interpolant matrix SOS cone, approximate Fekete nodes, no Gram variables\\
\code{sosineq\_interval}, \code{sosmatrixineq\_interval} & one & weighted cone on $[a,b]$, no multipliers, no Gram variables\\
\code{sosineq(prog,p,[a inf])} & one & weighted cone on a half-line\\
positivity on a box or semialgebraic set in several variables & any number & SOSTOOLS multipliers: \code{sossosvar\_nosdp} (scalar, Gram-free) or \code{sossosvar}, \code{sossosmatrixvar} (their Gram blocks become PSD cones)\\
\bottomrule
\end{tabular}
\end{center}
An interval constraint may depend only on its interval variable and on decision variables; the bridge stops with an error otherwise. Multivariate weighted cones with several weights (for example boxes without multipliers) are supported by the core (\code{nsd\_cone} accepts any weighted bases) but have no SOSTOOLS front end yet. The tests include matrix SOS constraints in two and three variables and box constraints in two variables (\cref{sec:verif}). Global constraints in several variables have been tested far less than univariate ones, and their cost grows quickly with the number of nodes $U=\binom{n+2d}{n}$.

%======================================================================
\section{Installation}\label{sec:install}
%======================================================================
\paragraph{Requirements.} MATLAB R2016b or later, or GNU Octave 8, and SOSTOOLS~4 (tested with the version on GitHub, \url{https://github.com/oxfordcontrol/SOSTOOLS}, October 2026). Programs must be written in SOSTOOLS' default polynomial mode (\code{pvar}). No SDP solver is needed.

\paragraph{MATLAB.}
\begin{lstlisting}
addpath(genpath('/path/to/SOSTOOLS'));
addpath('/path/to/sosnosdp');
setup_nosdp
\end{lstlisting}
\code{setup\_nosdp} adds the folders \code{core}, \code{bridge} and \code{examples} to the path. SOSTOOLS refuses to build a program unless an SDP solver is on the path. If none is installed, \code{setup\_nosdp} adds a placeholder \code{sedumi.m} that only raises an error when called. It is not added when a real solver is present.

\paragraph{GNU Octave.} SOSTOOLS~4 relies on three MATLAB behaviours that Octave lacks (sparse logical indexing and broadcasting of sparse arrays). The script \code{octave/patch\_sostools.py} applies three one-line edits to a copy of SOSTOOLS, which change nothing under MATLAB:
\begin{lstlisting}[language=bash]
git clone https://github.com/oxfordcontrol/SOSTOOLS.git
python3 sosnosdp/octave/patch_sostools.py SOSTOOLS
octave-cli --eval "SOSTOOLS_DIR='SOSTOOLS'; run('sosnosdp/octave/run_tests_octave.m')"
\end{lstlisting}
Two limitations of SOSTOOLS remain in Octave: products of a polynomial by a \code{dpvar} fail (write the \code{dpvar} first, for example \code{s*(x*(2-x))}), and the default form of \code{sosmatrixineq} fails. See \code{octave/README\_OCTAVE.md}.

%======================================================================
\section{Quick start}\label{sec:quick}
%======================================================================
Largest $a$ such that $M(x)=\begin{bmatrix}1+x^2&ax\\ax&2+x^2\end{bmatrix}$ is positive semidefinite for all $x$. The exact answer is $1+\sqrt2$.
\begin{lstlisting}
pvar x a

M = [1 + x^2, a*x;
     a*x,     2 + x^2];

prog = sosprogram(x, a);
prog = sosmatrixineq(prog, M);      % matrix SOS constraint
prog = sossetobj(prog, -a);         % maximize a

[prog, info] = sossolve_nosdp(prog);

double(sosgetsol(prog, a))          % 2.41421343
info.status                         % 'optimal'
\end{lstlisting}
The only change with respect to a SOSTOOLS script is the call to \code{sossolve\_nosdp}.

%======================================================================
\section{How it works}\label{sec:how}
%======================================================================
\subsection{SOS programs and the semidefinite route}
A symmetric $r\times r$ matrix polynomial $P$ of degree $2d$ is a sum of squares if $P=\sum_lQ_l^\top Q_l$ for polynomial matrices $Q_l$. With a basis $b$ of the polynomials of degree at most $d$ this holds if and only if $P(x)=(b(x)\otimes I_r)^\top S\,(b(x)\otimes I_r)$ for some $S\succeq0$~\cite{Nesterov2000,ShererHol2006}. Weighted SOS (WSOS) polynomials $\sum_jg_j\sigma_j$, with $\sigma_j$ SOS and weights $g_j\ge0$ on a set $D$, certify positivity on $D$~\cite{ShererHol2006}. An SOS program asks for decision variables $z$ such that affine polynomial expressions in $z$ are (W)SOS. SOSTOOLS introduces one Gram matrix $S$ per constraint and solves the resulting SDP.

\subsection{Interpolant representation and the matrix WSOS cone}
Following~\cite{LofbergParrilo2004,PappYildiz2019}, a polynomial of degree at most $2d$ is stored by its values at $N$ nodes $\xi_1,\dots,\xi_N$ that are unisolvent, meaning that only the zero polynomial vanishes at all of them. For an $r\times r$ matrix polynomial this is the vector $(\svec P(\xi_1),\dots,\svec P(\xi_N))$, where $\svec$ stacks the upper triangle with off-diagonal entries scaled by $\sqrt2$.

Let $\tilde p_{ji}=\sqrt{g_j(\xi_i)}\,b_j(\xi_i)$ be the weighted basis vectors at the nodes. The matrix WSOS cone and its dual are~\cite{KCV2023,PappAlizadeh2013}
\begin{gather*}
K=\Big\{\big(\textstyle\sum_j(\tilde p_{ji}\otimes I)^\top\Phi_j(\tilde p_{ji}\otimes I)\big)_{i=1}^N:\ \Phi_j\succeq0\Big\},\\
K^\ast=\Big\{W:\ \Lambda_j(W)=\textstyle\sum_i\tilde p_{ji}\tilde p_{ji}^\top\otimes W_i\succeq0\ \ \forall j\Big\}.
\end{gather*}
Membership in $K^\ast$ is a Cholesky test, and $F(W)=-\sum_j\log\det\Lambda_j(W)$ is a logarithmically homogeneous self-concordant barrier for $K^\ast$ with parameter $r\sum_jL_j$, where $L_j$ is the length of $b_j$. With $M^{(j)}_{ik}=(\tilde p_{ji}\otimes I)^\top\Lambda_j(W)^{-1}(\tilde p_{jk}\otimes I)$,
\[
\nabla F(W)_i=-\sum_jM^{(j)}_{ii},\qquad \nabla^2F(W)[H,H']=\sum_j\sum_{i,k}\tr\big(M^{(j)}_{ki}H_iM^{(j)}_{ik}H'_k\big).
\]
For $r=1$ these are the formulas of~\cite{PappYildiz2019}, and for general $r$ those of~\cite[Sec.~5.1]{KCV2023}. The PSD cone is the special case $N=1$, $\tilde p=1$. These formulas are implemented in \code{core/nsd\_cone\_eval.m}. The bases $\tilde P_j$ are orthonormalized, as recommended in~\cite[Sec.~4.1]{PappYildiz2019}.

\subsection{The interior-point method}
After the translation of \cref{sec:bridge}, every program has the form
\[
\max_z\ b^\top z\quad\text{s.t.}\quad c+Lz\in\cK=K_1\times\dots\times K_p ,
\]
which is the dual of $\min c^\top x$ s.t.\ $-L^\top x=b$, $x\in\cK^\ast$. The homogeneous self-dual embedding of this pair is solved by the predictor-corrector method of Skajaa and Ye~\cite{SkajaaYe2015}, as in~\cite{PappYildiz2019,PappYildiz2017}. It only needs the barrier $F$ of $\cK^\ast$, and the SOS cone $\cK$ appears only through the dual slack. The Newton systems are solved through a Schur complement of size equal to the number of decision variables (\code{core/nsd\_ipm.m}). The implementation uses practical neighbourhood parameters (predictor to the boundary of a neighbourhood of size 0.7, up to six correctors back to 0.3) rather than the theoretical ones of~\cite{PappYildiz2017}, so the worst-case iteration bound is not claimed. Correctness rests on the certificates of \cref{sec:cert}.

\subsection{Certificates}\label{sec:cert}
The method approaches $\cK$ from outside, through the dual iterates, and $\cK$ has no direct membership test. Every returned solution is therefore certified (\code{core/nsd\_certify\_point.m}, \code{core/nsd\_check\_values.m}).
\begin{enumerate}[itemsep=2pt]
\item \emph{Projection.} At an iterate with $\delta=\mu/\tau$, the decision variables are re-solved as $z^\star=\arg\min_z\|c+Lz+\delta\nabla F(x)\|_{\nabla^2F(x)^{-1}}$, so that $S=c+Lz^\star$ lies as close as possible to the central point $\hat s=-\delta\nabla F(x)$.
\item \emph{Gram test} (decides in the default mode). The Gram matrices $\Phi_j=\Lambda_j(x)^{-1}\Lambda_j(w)\Lambda_j(x)^{-1}$, $w=\nabla^2F(x)^{-1}S$, are formed. This is the recovery of~\cite[Thm.~5.1]{PappYildiz2019} in matrix form: $\sum_j\cA_j(\Phi_j)=S$, and if $\|S+\delta\nabla F(x)\|_{\nabla^2F(x)^{-1}}<\delta$ all $\Phi_j$ are positive definite. Because $\sum_j\cA_j(\Lambda_j^{-1}\Lambda_j(w)\Lambda_j^{-1})=\nabla^2F(x)w$, the same formula applied to the residual gives a correction, and the identity is refined iteratively to rounding level. The test accepts if every $\Phi_j$ satisfies $\lambda_{\min}>10^{-12}\lambda_{\max}$ and the relative residual of the identity is below $10^{-10}$ and below $10^{-2}\lambda_{\min}/\lambda_{\max}$, so that the residual cannot destroy positive definiteness. The Gram matrices are used only here, never in the solve.
\item \emph{Dikin test} (reported, decides only in mode \code{'dikin'}). If $r=\|S-\hat s\|_{\nabla^2F(x)^{-1}}/\delta<1$, then $S$ lies in the Dikin ellipsoid of the conjugate barrier at $\hat s$, hence in the interior of $\cK$, with the scale-free depth $(1-r)\,\delta\,\|E\|/(\lambda_{\max}(\Lambda(x))\|S\|)$ for the interior element $E=\cA(I)$. This test needs no Gram matrix, but it relies on the inverse Hessian, which is ill-conditioned near the boundary. In floating point it accepted a matrix polynomial that is not positive semidefinite far from the nodes (test T15), with a computed depth of $10^{-8}$, so it is not used to decide by default.
\item \emph{Returned coefficients.} The certificate is evaluated for the values $c+Lz$ of the coefficients actually returned to \code{sosgetsol}. The last ten certified iterates are kept, and the newest one that passes this final check is returned.
\end{enumerate}
For infeasible programs the method returns $x\in\cK^\ast$ with $c^\top x<0$. Then no $z$ with $\|z\|<|c^\top x|/\|L^\top x\|$ is feasible (\code{info.farkas\_radius}). This is a global proof of infeasibility only when the radius is large compared with the expected size of solutions.

\subsection{From a SOSTOOLS program to cones}\label{sec:bridge}
SOSTOOLS stores every constraint as polynomial coefficients that are affine in the decision vector $u$, namely $b-A^\top u$ over a list of monomials (fields \code{prog.expr.At}, \code{b}, \code{Z}). The Gram matrices are created only inside \code{sossolve}. \code{sossolve\_nosdp} reads the program at that point (\code{bridge/sossolve\_nosdp.m}):
\begin{enumerate}[itemsep=2pt]
\item \emph{Decision variables.} Free variables (\code{sosdecvar}, \code{sospolyvar}, \code{sospolymatrixvar}) are kept. Gram blocks that SOSTOOLS creates for \code{sossosvar}, \code{sossosmatrixvar} and \code{sosposmatr} are parametrized by their symmetric part and constrained to the PSD cone of the same method. \code{sossosvar\_nosdp} avoids these Gram blocks altogether.
\item \emph{Global SOS constraints} (\code{sosineq}, \code{sosmatrixineq}), in any number of variables. The monomials of the Gram basis are taken from the Newton box of the polynomial. Nodes unisolvent for their products are approximate Fekete points, selected by column-pivoted QR from Chebyshev points (one variable) or Halton points (several variables)~\cite{SommarivaVianello2009,PappYildiz2019}. Monomials that cannot occur in an SOS polynomial become equality constraints. The default form of \code{sosmatrixineq}, $y^\top M(x)y$ with auxiliary variables $y$, is converted back to the matrix $M(x)$, so that the matrix cone of~\cite{KCV2023} is used rather than a scalar cone in $(x,y)$.
\item \emph{Interval constraints} (one variable). For \code{sosineq\_interval} and \code{sosmatrixineq\_interval} on $[a,b]$ the cone has weights $1$ and $(t-a)(b-t)$, Chebyshev--Lobatto nodes and orthonormalized Chebyshev bases. Half-line constraints of SOSTOOLS (\code{sosineq(prog,p,[a inf])}) use weights $1$ and $t-a$. No multiplier is introduced.
\item \emph{Equalities} (\code{soseq}) are eliminated through a null-space basis.
\item \emph{Reparametrization.} The constraint matrix is factored as $L\Pi=QR$ with column pivoting, and the solver works with the orthonormal coordinates $Q$. The problem depends only on the range of $L$, so this is exact. It makes the method independent of the user's basis and scaling, and removes redundant variables. Numerically dependent directions are reported (\code{info.near\_dependent}).
\item \emph{Write-back.} The SOSTOOLS decision vector is reconstructed and stored in the field \code{prog.solinfo.RRx}, which is what \code{sosgetsol} reads.
\end{enumerate}

\subsection{Size and cost compared with the SDP}\label{sec:cost}
Consider one $r\times r$ constraint of degree $2d$ in one variable with weights $1$ and $(t-a)(b-t)$. The cone has $U=2d+1$ nodes and dimension $m=Ur(r+1)/2=\Theta(r^2d)$, and the Gram blocks of the SDP have sides $n_s=r(d+1)$ and $rd$. The SDP has $m$ equality constraints and $\Theta(r^2d^2)$ Gram unknowns, the interpolant formulation has no Gram unknowns. Both have the barrier parameter $r(2d+1)$, hence the same iteration bound $O(\sqrt{rd}\,\log\varepsilon^{-1})$~\cite{SkajaaYe2015,PappYildiz2019}. With dense linear algebra, one iteration costs:
\begin{center}\small
\begin{tabular}{lll}
\toprule
& operations per iteration & memory\\
\midrule
this method: barrier gradient and Hessian & $O(r^4d^3)$ & $O(r^4d^2)$\\
this method: Newton system (order $m$) & $O(m^3)=O(r^6d^3)$ & $O(r^4d^2)$\\[2pt]
SDP, Schur complement: forming it & $O(mn_s^3+m^2n_s^2)=O(r^6d^4)$ & $O(r^4d^2)$\\
SDP, Schur complement: factoring it & $O(m^3)=O(r^6d^3)$ & \\[2pt]
SDP, full KKT in all Gram variables & $O(r^6d^6)$ & $O(r^4d^4)$\\
\bottomrule
\end{tabular}
\end{center}
Schur-complement solvers include SeDuMi, SDPT3 and MOSEK. The dominant term of this method, the factorization of a dense matrix of order $m$, has the same order as the factorization of the Schur complement in an SDP solver. The difference per iteration is the formation of the Schur complement, larger by a factor of order $L^2/U\approx d/2$ when the constraint matrices are treated as dense. Solvers that exploit their sparsity or low rank can reduce this term, for instance with the rank-one constraint matrices of an interpolant formulation~\cite{LofbergParrilo2004}. The bound $O(L^6)$ per iteration for the SDP in~\cite[Sec.~1.2]{PappYildiz2019}, here $O(r^6d^6)$, counts the solution of a linear system in all $\Theta(L^2)$ Gram variables, which is the last row of the table. The time per iteration of MOSEK reported in~\cite[Table~1]{PappYildiz2019} grows roughly like $d^4$ (from 0.92~s at $d=100$ to 661~s at $d=500$).

In $n$ variables, with $L=\binom{n+d}{n}$ and $U=\binom{n+2d}{n}$, the factorization costs $O(r^6U^3)$ in both formulations and the formation of the Schur complement $O(r^6U^2L^2)$. The ratio $L^2/U$ is about $L/2^n$ for fixed $n$, so it grows with the degree, and it is bounded when $d$ is fixed and $n$ grows, in which case both formulations cost $O(r^6L^6)$ per iteration.

In short, against a well-implemented Schur-complement SDP solver the gain per iteration is at most of order $L^2/U$, and it can vanish. The practical advantages of the method are the absence of Gram variables, a memory footprint independent of the Gram blocks, the conditioning of the interpolant representation~\cite[Sec.~4]{PappYildiz2019}, and lower constants. Against solvers that factor the full KKT system the gain is of order $d^3$ in time and $d^2$ in memory in one variable.

%======================================================================
\section{Reference}\label{sec:ref}
%======================================================================
\subsection{Functions}
\begin{description}[style=nextline,itemsep=3pt]
\item[\code{[prog, info] = sossolve\_nosdp(prog, opts)}] Solves a SOSTOOLS program without SDP. Options and outputs are listed below.
\item[\code{prog = sosineq\_interval(prog, p, t, [a b])}] $p(t)\ge0$ for $t\in[a,b]$. $p$ may depend on decision variables and on the program variable $t$ only.
\item[\code{prog = sosmatrixineq\_interval(prog, M, t, [a b])}] $M(t)\succeq0$ for $t\in[a,b]$, $M$ square and symmetric.
\item[\code{[prog, s] = sossosvar\_nosdp(prog, Z)}] SOS decision polynomial in the monomial basis $Z$, represented by its coefficients and its own SOS cone. The program remains solvable by \code{sossolve}.
\item[\code{[s, tm] = nsd\_timevar('s', [a b])}] Normalized variable $s\in[-1,1]$ for a physical interval, with \code{tm.ddt(Y)} ($dY/dt$), \code{tm.at(Y, t0)} ($Y$ at time $t_0$), \code{tm.interval(I)} (an interval in $s$) and \code{tm.t(s0)}.
\item[\code{setup\_nosdp}] Adds the package to the path.
\item[\code{nsd\_cone}, \code{nsd\_solve}] Low-level interface: \code{K = nsd\_cone(r, Ps)} builds a cone from the weighted basis matrices, and \code{out = nsd\_solve(cones, L, c, b, opts)} solves $\max b^\top z$ s.t.\ $c+Lz\in\cK$ without SOSTOOLS.
\end{description}

\subsection{Options of \texttt{sossolve\_nosdp}}
\begin{center}\small
\begin{tabular}{lll}
\toprule
Option & Default & Meaning\\
\midrule
\code{certificate} & \code{'auto'} & \code{'auto'} and \code{'gram'}: the Gram test decides; \code{'dikin'}: Gram-free\\
\code{depth\_tol} & $10^{-8}$ & depth a Dikin certificate must prove (mode \code{'dikin'})\\
\code{max\_iter}, \code{tol} & 200, $10^{-8}$ & iterations and tolerance of the interior-point method\\
\code{box} & \code{[-1 1]} & box in which nodes of global SOS constraints are chosen\\
\code{rank\_tol} & $10^{-14}$ & tolerance for removing dependent decision directions\\
\code{verbose} & \code{false} & print the iterations\\
\bottomrule
\end{tabular}
\end{center}

\subsection{Statuses and output}
\begin{description}[itemsep=2pt]
\item[\code{feasible}] feasibility program, certified solution.
\item[\code{optimal}] the method converged; the returned solution is the newest certified iterate that passes the final check, strictly feasible and slightly inside the optimum (for SOS programs the optimum lies on the boundary of the cone). In the tests the gap is between $10^{-9}$ and $10^{-6}$ relative.
\item[\code{infeasible}] a Farkas vector was found; see \code{info.farkas\_radius}.
\item[\code{uncertified\_*}] no certificate could be established; the last iterate is still returned so that it can be checked by other means. \code{uncertified\_coefficients} means that no certified iterate survived the check on the returned coefficients.
\end{description}
Further fields of \code{info}: \code{iters}, \code{time}, \code{obj}, \code{cones} (one description per cone), \code{cert\_method} and \code{depth} (per cone), \code{dikin\_ratio}, \code{rank}, \code{near\_dependent}, \code{barrier\_parameter}, \code{accuracy} (\code{full} or \code{reduced}), \code{last\_failed\_certificate}.

%======================================================================
\section{Good practice and numerical limits}\label{sec:practice}
%======================================================================
\paragraph{Normalize interval variables.} SOSTOOLS stores all polynomials in monomial coefficients of the declared variable. On a wide interval high-degree monomials are numerically dependent: for a clock-dependent LMI on $[0,2]$ the numerical rank drops from 249 to 92 at degree 40 and the method stalls, whereas in a variable normalized to $[-1,1]$ it certifies degree 40 in 19 iterations. The change of basis cannot be repaired after the program is built, so declare the variable with \code{nsd\_timevar}, which still lets you write derivatives and endpoint values in physical time.

\paragraph{Prefer interval constraints to multipliers.} \code{sosmatrixineq\_interval} needs neither multipliers nor Gram variables. On a clock-dependent LMI of degree 24 it took 0.8~s, whereas the same solver with a matrix SOS multiplier did not reach a certificate in 21~s.

\paragraph{Constraints on the boundary of the SOS cone.} Certificates prove membership in the \emph{interior} of the cone. Some constraints lie on its boundary for every value of the decision variables, and then no certificate can be given. A typical cause is an entry of lower degree than the diagonal pattern requires. For example, in $\begin{bmatrix}1+x^2+y^2&a(x+y)\\a(x+y)&2\end{bmatrix}$ the constant $(2,2)$ entry forces the Gram blocks of $x$ and $y$ in that entry to vanish, so every Gram matrix is singular (test T15). The status is then \code{uncertified\_*}. Remedies are to give every diagonal entry the full degree, if the problem allows it, or to reduce the monomial basis of the constraint.

\paragraph{Ill-conditioned problems.} When every solution is badly conditioned, no floating-point certificate can be trusted, whatever the solver. For a random 10-state plant with one input and sampling period 0.1, the discrete controllability Gramian has condition number about $10^{17}$ and every valid Lyapunov matrix about $10^{7}$. The status is then \code{uncertified\_stalled}, while the returned controller still stabilizes the sampled loop (spectral radius 0.98). For control problems, an independent check of the result (closed-loop spectral radius, second moments) is always recommended.

%======================================================================
\section{Examples}\label{sec:examples}
%======================================================================
The examples in \code{examples/} are written one statement per line with comments, so that they can serve as templates.

\subsection{Example 1: polynomial minimization}
Global minimum of $x^4-4x^3+3x^2+2x+1$ (0.75) and minimum of $x^4-x$ on $[0,2]$, the latter imposed directly on the interval:
\begin{lstlisting}
prog = sosprogram(x, gam);
prog = sosineq_interval(prog, x^4 - x - gam, x, [0 2]);
prog = sossetobj(prog, -gam);
[prog, info] = sossolve_nosdp(prog);
\end{lstlisting}
Result: $-0.4724704$ (exact $-0.4724704$), certified.

\subsection{Example 2: matrix SOS}
The quick-start problem, in both forms of \code{sosmatrixineq}. Result $2.4142134$ (exact $1+\sqrt2$).

\subsection{Example 3: sampled-data stabilization by clock-dependent LMIs}
For the plant $\dot x=Ax+Bu$ with a zero-order hold and sampling period $T$, the augmented state $(x,u)$ yields an impulsive system~\cite{Briat2013,Briat2021}. One looks for $X(t)$, $t\in[0,T]$, and $U$ with $\dot X-A_aX-XA_a^\top\succeq0$ on $[0,T]$ and a boundary LMI, and $K=UX(T)^{-1}$ stabilizes the sampled loop.
\begin{lstlisting}
[s, tm] = nsd_timevar('s', [0 T]);
prog = sosprogram(s);

[prog, X] = sospolymatrixvar(prog, monomials(s, 0:d), [m m], 'symmetric');
[prog, U] = sospolymatrixvar(prog, monomials(s, 0), [nu m]);

% flow condition on [0, T]
P = tm.ddt(X) - X*Aa' - (X*Aa')';
prog = sosmatrixineq_interval(prog, P, s, tm.interval());

% jump condition
XT = tm.at(X, T);
X0 = tm.at(X, 0);
Off = -(XT*J0' + U'*Ba')';
prog = sosmatrixineq(prog, [XT, Off'; Off, X0 - epsl*eye(m)], 'Mineq');

[prog, info] = sossolve_nosdp(prog);
\end{lstlisting}
For the unstable oscillator $A=\begin{bmatrix}0.3&3\\-3&0.3\end{bmatrix}$, $B=[0\ 1]^\top$, $T=2$: infeasible at $d=6$, feasible at $d=8$ (spectral radius 0.904) and at $d=32$ (0.539).

\subsection{Example 4: largest sampling interval under aperiodic sampling}
With a given gain $K$ and sampling intervals in $[T_{\min},T_{\max}]$, the jump condition must hold for every interval $\theta\in[T_{\min},T_{\max}]$: it is a parameter-dependent LMI, imposed with \code{sosmatrixineq\_interval(prog, Psi, s, tm.interval([Tmin Tmax]))}. The largest certified $T_{\max}$ is found by bisection. For the benchmark $A=\begin{bmatrix}0&1\\0&-0.1\end{bmatrix}$, $B=[0\ 0.1]^\top$, $K=-[3.75\ 11.5]$, $T_{\min}=0.01$ and degree 8 the result is $T_{\max}=1.7292$, against the periodic bound $1.7298$ computed on the same grid.

\subsection{Example 5: stochastic sampled-data systems}
Theorem~15 of~\cite{Briat2016} gives conditions for the mean-square stabilization of
\[
dx=(A_{sd}x+B_{sd}u)\,dt+E_{sd}x\,dW_1+\alpha B_{sd}u\,dW_2,\qquad u(t)=K_d\,[x(t_k);u(t_k)],
\]
under sampling intervals in $[T_{\min},T_{\max}]$: a $9\times9$ matrix inequality on $[0,T_{\max}]$ and a $6\times6$ one on $[T_{\min},T_{\max}]$, both imposed directly on their intervals (\code{examples/stochastic\_sd\_design.m}). For the paper's Example~16 all six ranges of its Table~6 are certified in under 0.5~s each. Each design is verified exactly through the second moment $E[xx^\top]$, whose evolution between and at samples is linear (\code{examples/second\_moment.m}). For an unstable variant, bisection gives the largest certified $T_{\max}$: 1.37 at degree 2 and 3.47 at degree 4.

%======================================================================
\section{Verification}\label{sec:verif}
%======================================================================
\code{tests/run\_tests.m} contains 21 tests with known answers. They cover polynomial minima, the matrix example, half-line and interval minima, multipliers through \code{sossosvar} and \code{sossosvar\_nosdp}, equality constraints, clock-dependent LMIs that are infeasible at degree 6 and feasible at degrees 8, 24, 32 and 40, Gram-free certificates, \code{nsd\_timevar}, matrix SOS in two variables and scalar SOS in three variables with exact answers, scalar and matrix inequalities on a box in two variables through multipliers, and two regression tests ensuring that points on the boundary of the cone are not accepted as certificates (T12, T15). All pass under Octave 8.4 with SOSTOOLS~4. The workflow \code{.github/workflows/octave-tests.yml} runs them on every push.

With SeDuMi through SOSTOOLS' \code{sossolve} (\code{tests/compare\_sedumi.m}), the optimal values of four test problems agree within $1.4\cdot10^{-6}$, and within $2\cdot10^{-7}$ for three of them. The values of SOSNOSDP lie slightly inside the feasible set, those of SeDuMi sometimes slightly outside. On the clock-dependent LMI of Example~3 at increasing degree (\code{tests/scaling\_sedumi.m}, time interval $[0,2]$):
\begin{center}\small
\begin{tabular}{rlll}
\toprule
degree & SeDuMi, multiplier form & SOSNOSDP, multiplier form & SOSNOSDP, native interval\\
\midrule
8 & feasible, 0.6 s & feasible, 0.7 s & feasible, 0.4 s\\
16 & numerical error flag & feasible, 2.8 s & feasible, 0.5 s\\
24 & numerical error & uncertified & feasible, 0.8 s\\
\bottomrule
\end{tabular}
\end{center}
These are single runs on one core under Octave.

In several variables (\code{tests/run\_tests.m}, T13--T17, and SeDuMi through \code{sossolve}; the exact value of the first problem is $1+1/\sqrt2$, and the values of SOSNOSDP are certified):
\begin{center}\small
\begin{tabular}{llll}
\toprule
problem & exact & SeDuMi & SOSNOSDP\\
\midrule
$2\times2$ matrix SOS, 2 variables, degree 2 & $1.7071068$ & 1.7071068 & 1.7071067\\
$3\times3$ matrix SOS, 3 variables, degree 2 & & 1.7320508 & 1.7320508\\
scalar SOS, 3 variables, degree 4 & 4 & 4.0000000 & 3.9999997\\
$\min xy$ on $[-1,1]^2$, Gram-free multipliers & $-1$ & & $-1.00000002$\\
$2\times2$ matrix on $[-1,1]^2$, multipliers & 2 & 2.0000000 & 1.9999999\\
\bottomrule
\end{tabular}
\end{center}
The package has not yet been run in MATLAB by its authors.

%======================================================================
\section{Relation to other software}\label{sec:related}
%======================================================================
\begin{itemize}[itemsep=2pt]
\item \textbf{alfonso}~\cite{alfonso} (MATLAB, Papp and Y{\i}ld{\i}z): generic non-symmetric conic solver with scalar WSOS cones in the interpolant basis, and the reference implementation of~\cite{PappYildiz2019}. It accepts user-defined barrier oracles. \url{https://github.com/dpapp-github/alfonso}
\item \textbf{Hypatia.jl}~\cite{CKV2022,CKV2023} (Julia, Coey, Kapelevich and Vielma): generic conic interior-point solver with, among many cones, the scalar and matrix interpolant WSOS cones (\code{WSOSInterpNonnegativeCone}, \code{WSOSInterpPosSemidefTriCone}) of~\cite{KCV2023}. \url{https://github.com/jump-dev/Hypatia.jl}
\item \textbf{SumOfSquares.jl}~\cite{SumOfSquares} (Julia, Legat et al.): JuMP-based modelling of SOS programs, the Julia counterpart of SOSTOOLS. With an interpolant (Lagrange) basis, scalar SOS constraints reach Hypatia's interpolant cone without a PSD reformulation. Matrix SOS constraints are scalarized. \url{https://github.com/jump-dev/SumOfSquares.jl}
\item \textbf{SOSTOOLS}~\cite{SOSTOOLS} (MATLAB): the modelling front end used by SOSNOSDP, which normally solves programs through SDP solvers. \url{https://github.com/oxfordcontrol/SOSTOOLS}
\end{itemize}
SOSNOSDP is not a port of Hypatia. Its solver is a compact MATLAB implementation of the same cone and of a Skajaa--Ye type method, and its purpose is to make the approach available to existing SOSTOOLS programs, in particular matrix-valued ones with interval constraints.

%======================================================================
\section{Limitations and future work}\label{sec:limits}
%======================================================================
\begin{itemize}[itemsep=2pt]
\item Only SOSTOOLS programs written with \code{pvar} are supported, not the Symbolic Math Toolbox mode.
\item The sparsity options of SOSTOOLS (\code{'sparse'}, \code{'sparsemultipartite'}) are handled as dense constraints.
\item Global SOS constraints in several variables use approximate Fekete nodes from Halton points and have been tested far less than univariate and interval constraints.
\item Linear algebra is dense. The Hessian of a cone with $N$ nodes and matrix size $r$ has order $Nr(r+1)/2$.
\item Interval constraints are univariate. Boxes in several variables need SOSTOOLS multipliers; a Gram-free front end for multivariate weighted cones (robust or LPV problems) is a natural extension, since the core already supports them. There is no Gram-free matrix multiplier yet (\code{sossosmatrixvar} uses PSD Gram blocks).
\item Constraints lying on the boundary of the SOS cone for every value of the decision variables cannot be certified (\cref{sec:practice}).
\item The interior-point method uses practical parameters and has not been tuned. Certificates are evaluated in floating point, not in exact arithmetic.
\end{itemize}

%======================================================================
\section{How to cite}\label{sec:cite}
%======================================================================
If you use SOSNOSDP, please cite the package (\code{CITATION.cff}) together with the works it builds on: Papp and Y{\i}ld{\i}z~\cite{PappYildiz2019} for SOS optimization without SDP, Kapelevich, Coey and Vielma~\cite{KCV2023} for the matrix SOS cone and its barrier, Coey, Kapelevich and Vielma~\cite{CKV2022} for Hypatia, Skajaa and Ye~\cite{SkajaaYe2015} for the interior-point method, and SOSTOOLS~\cite{SOSTOOLS}.

\begin{thebibliography}{99}\small
\bibitem{Briat2013} C.~Briat, \emph{Convex conditions for robust stability analysis and stabilization of linear aperiodic impulsive and sampled-data systems under dwell-time constraints}, Automatica, 49 (2013), pp.~3449--3457.
\bibitem{Briat2016} C.~Briat, \emph{Stability analysis and stabilization of stochastic linear impulsive, switched and sampled-data systems under dwell-time constraints}, Automatica, 74 (2016), pp.~279--287.
\bibitem{Briat2021} C.~Briat, \emph{Stability analysis and stabilization of LPV systems with jumps and (piecewise) differentiable parameters using continuous and sampled-data controllers}, Nonlinear Analysis: Hybrid Systems, 41 (2021), 101040.
\bibitem{CKV2022} C.~Coey, L.~Kapelevich, and J.~P.~Vielma, \emph{Solving natural conic formulations with Hypatia.jl}, INFORMS Journal on Computing, 34 (2022), pp.~2686--2699.
\bibitem{CKV2023} C.~Coey, L.~Kapelevich, and J.~P.~Vielma, \emph{Performance enhancements for a generic conic interior point algorithm}, arXiv:2107.04262.
\bibitem{KCV2023} L.~Kapelevich, C.~Coey, and J.~P.~Vielma, \emph{Sum of squares generalizations for conic sets}, Mathematical Programming (2023), pp.~1417--1429, \href{https://doi.org/10.1007/s10107-022-01831-6}{doi:10.1007/s10107-022-01831-6}. Preprint arXiv:2103.11499.
\bibitem{LofbergParrilo2004} J.~L\"ofberg and P.~A.~Parrilo, \emph{From coefficients to samples: a new approach to SOS optimization}, in 43rd IEEE Conference on Decision and Control, 2004, pp.~3154--3159.
\bibitem{Nesterov2000} Y.~Nesterov, \emph{Squared functional systems and optimization problems}, in High Performance Optimization, Kluwer, 2000, pp.~405--440.
\bibitem{PappAlizadeh2013} D.~Papp and F.~Alizadeh, \emph{Semidefinite characterization of sum-of-squares cones in algebras}, SIAM Journal on Optimization, 23 (2013), pp.~1398--1423.
\bibitem{PappYildiz2017} D.~Papp and S.~Y{\i}ld{\i}z, \emph{On ``A homogeneous interior-point algorithm for non-symmetric convex conic optimization''}, arXiv:1712.00492, 2017.
\bibitem{PappYildiz2019} D.~Papp and S.~Y{\i}ld{\i}z, \emph{Sum-of-squares optimization without semidefinite programming}, SIAM Journal on Optimization, 29 (2019), pp.~822--851.
\bibitem{alfonso} D.~Papp and S.~Y{\i}ld{\i}z, \emph{alfonso: Matlab package for nonsymmetric conic optimization}, INFORMS Journal on Computing, 34 (2022), pp.~11--19.
\bibitem{ShererHol2006} C.~W.~Scherer and C.~W.~J.~Hol, \emph{Matrix sum-of-squares relaxations for robust semi-definite programs}, Mathematical Programming, 107 (2006), pp.~189--211.
\bibitem{SkajaaYe2015} A.~Skajaa and Y.~Ye, \emph{A homogeneous interior-point algorithm for nonsymmetric convex conic optimization}, Mathematical Programming, 150 (2015), pp.~391--422.
\bibitem{SommarivaVianello2009} A.~Sommariva and M.~Vianello, \emph{Computing approximate Fekete points by QR factorizations of Vandermonde matrices}, Computers \& Mathematics with Applications, 57 (2009), pp.~1324--1336.
\bibitem{SOSTOOLS} A.~Papachristodoulou, J.~Anderson, G.~Valmorbida, S.~Prajna, P.~Seiler, P.~A.~Parrilo, M.~M.~Peet, and D.~Jagt, \emph{SOSTOOLS: Sum of squares optimization toolbox for MATLAB}, version 4.00, 2021. \url{https://github.com/oxfordcontrol/SOSTOOLS}.
\bibitem{Sturm1999} J.~F.~Sturm, \emph{Using SeDuMi 1.02, a MATLAB toolbox for optimization over symmetric cones}, Optimization Methods and Software, 11--12 (1999), pp.~625--653.
\bibitem{SumOfSquares} B.~Legat, C.~Coey, R.~Deits, J.~Huchette, and A.~Perry, \emph{Sum-of-squares optimization in Julia}, 2017. \url{https://github.com/jump-dev/SumOfSquares.jl}.
\end{thebibliography}

\end{document}
